\editionPage[\pdfbookmark[0]{Démonstrations}{bm-preuves-debut-parent}]{preuves-debut}{Comprendre les preuves}{Pourquoi le produit de deux négatifs est-il positif ?}{Les arguments ci-dessous éclairent des propriétés utilisées dans le livre.}
\begin{blocEdition}{preuves-debut}

\editionEtape{Partir de l’opposé}Soient $a$ et $b$ positifs. Comme $b+(-b)=0$, multiplier par $a$ donne $ab+a(-b)=0$ par distributivité. Ainsi $a(-b)$ est l’opposé de $ab$ : $a(-b)=-ab$.
\editionEtape{Changer aussi le premier signe}De même, $a+(-a)=0$. En multipliant par $-b$, on obtient $a(-b)+(-a)(-b)=0$. On sait que $a(-b)=-ab$ ; donc $(-a)(-b)$ est l’opposé de $-ab$, soit $ab$.
\editionEtape{Ce que prouve le raisonnement}La règle des signes est cohérente avec la distributivité et la notion d’opposé. Elle ne repose pas seulement sur une convention isolée à mémoriser. Le cas d’un facteur nul se traite à part : son produit vaut 0.
\editionRetenir{Deux facteurs négatifs donnent un produit positif parce que ce produit doit être l’opposé d’un produit négatif pour que la distributivité reste vraie.}
\end{blocEdition}
\editionPage{preuves-2}{Comprendre les preuves}{Pourquoi diviser par une fraction revient-il à inverser ?}{Les arguments ci-dessous éclairent des propriétés utilisées dans le livre.}
\begin{blocEdition}{preuves-2}

\editionEtape{Utiliser la définition du quotient}Pour $v\ne0$, le quotient $u/v$ est l’unique nombre $q$ tel que $qv=u$.
\editionEtape{Tester une expression}L’inverse de $v$ vérifie $v\times(1/v)=1$. Le nombre $u\times(1/v)$, multiplié par $v$, donne $u\times[(1/v)\times v]=u$. C’est donc le quotient $u/v$.
\editionEtape{Appliquer aux fractions}Si $v=c/d$ avec $c,d\ne0$, son inverse est $d/c$ car $(c/d)(d/c)=1$. Ainsi $\frac ab:\frac cd=\frac ab\times\frac dc$, avec tous les dénominateurs non nuls.
\editionEtape{Un exemple contrôlé}$\frac34:\frac25=\frac34\times\frac52=\frac{15}{8}$. Vérifier en sens inverse : $\frac{15}{8}\times\frac25=\frac34$. On inverse uniquement le diviseur.
\editionRetenir{Pour justifier un quotient, remultiplier la réponse par le diviseur : on doit retrouver le dividende.}
\end{blocEdition}
\editionPage{preuves-3}{Comprendre les preuves}{Pythagore par une comparaison d’aires}{Les arguments ci-dessous éclairent des propriétés utilisées dans le livre.}
\begin{blocEdition}{preuves-3}

\editionEtape{Construire la figure}On place quatre triangles rectangles identiques, de côtés perpendiculaires $a,b$ et d’hypoténuse $c$, dans un carré de côté $a+b$. Le quadrilatère central a quatre côtés de longueur $c$. Ses angles sont droits car les deux angles aigus d’un triangle rectangle ont pour somme $90^\circ$ : c’est un carré.
\label{fig-preuves-3-preuve-pythagore}\figureGeo{preuve-pythagore}

\editionEtape{Compter deux fois la même aire}Le grand carré a pour aire $(a+b)^2$. La somme des aires des quatre triangles et du carré central est $4(ab/2)+c^2=2ab+c^2$.
\editionEtape{Comparer}En partageant un carré de côté $a+b$ en deux carrés et deux rectangles, on obtient aussi son aire $a^2+2ab+b^2$. Les deux calculs donnent $a^2+2ab+b^2=2ab+c^2$. En soustrayant $2ab$, il reste $\res{a^2+b^2=c^2}$.
\editionRetenir{Le théorème compare les aires des carrés construits sur les côtés. Cette preuve concerne le triangle rectangle, hypothèse essentielle.}
\end{blocEdition}
\editionPage{preuves-4}{Comprendre les preuves}{Pourquoi le cosinus dépend-il de l’angle ?}{Les arguments ci-dessous éclairent des propriétés utilisées dans le livre.}
\begin{blocEdition}{preuves-4}

\editionEtape{Deux triangles rectangles emboîtés}On considère un triangle ABC rectangle en B et un point M sur AB. La parallèle à BC par M coupe AC en N. Le triangle AMN est rectangle en M ; il a le même angle aigu en A que le grand triangle.
\label{fig-preuves-4-cours-thales}\figureGeo{cours-thales}

\editionEtape{Appliquer Thalès}On a $AM/AB=AN/AC$. Les longueurs étant non nulles, les produits en croix donnent $AM\times AC=AN\times AB$, puis $AM/AN=AB/AC$.
\editionEtape{Interpréter}Le rapport côté adjacent à A sur hypoténuse est donc le même dans les deux triangles. On peut agrandir ou réduire le triangle sans le changer tant que l’angle est conservé : c’est ce rapport qui définit le cosinus de l’angle.
\editionRetenir{Un cosinus est un rapport de longueurs, donc un nombre sans unité. Le choix de l’angle fixe le côté adjacent ; l’hypoténuse reste opposée à l’angle droit.}
\end{blocEdition}
